\documentclass[12pt,a4paper]{article}
\usepackage[utf8]{inputenc}
\usepackage{amsmath,amssymb,amsthm}
\usepackage{geometry}
\geometry{margin=2.5cm}
\usepackage{booktabs}
\usepackage{hyperref}
\usepackage{xcolor}

\hypersetup{
    colorlinks=true,
    linkcolor=blue!60!black,
    citecolor=blue!60!black,
    urlcolor=blue!60!black
}

\newcommand{\keV}{\,\mathrm{keV}}
\newcommand{\GeV}{\,\mathrm{GeV}}
\newcommand{\eV}{\,\mathrm{eV}}
\newcommand{\Mpc}{\,\mathrm{Mpc}}
\newcommand{\hMpc}{\,h\,\mathrm{Mpc}^{-1}}
\newcommand{\CP}{\mathbb{CP}}
\newcommand{\Sph}{S}

\newtheorem{theorem}{Theorem}
\newtheorem{lemma}[theorem]{Lemma}
\newtheorem{proposition}[theorem]{Proposition}
\newtheorem{corollary}[theorem]{Corollary}
\newtheorem{definition}[theorem]{Definition}
\newtheorem{remark}{Remark}

\title{The $\psi$ Field Alone Cannot Reproduce the CDM Transfer Function:\\
A Five-Argument Proof with Topological Corollary}
\author{DFD Research Programme}
\date{April 2026}

\begin{document}
\maketitle

\begin{abstract}
We prove that in Density Field Dynamics (DFD) with internal manifold
$\CP^2 \times \Sph^3$, the scalar refractive field $\psi$ arising from
the zeroth Betti number $b_0 = 1$ cannot, by itself, reproduce the
cold dark matter (CDM) transfer function $T(k)$ or the observed matter
power spectrum $P(k)$. We establish this through five independent
arguments: (1) the $\psi$ field equation is elliptic in the quasi-static
limit, making $\psi$ a constraint variable without independent dynamics;
(2) the DFD enhanced potential is an odd function of the Newtonian
potential, so its time-average over baryon acoustic oscillations
vanishes identically; (3) the temporal $\psi$ sector has a dust-like
equation of state but negligible amplitude $\Omega_\psi < 10^{-19}$;
(4) the maximum quantitative enhancement from $\psi$ is $\sim 30\times$,
whereas CDM provides $\sim 175\times$ at the scales controlling
$\sigma_8$; and (5) the $\psi$-screen effect is purely optical and
provides at most $\sim 2\times$ distance enhancement, orders of
magnitude below what is required. As a corollary, we show that the
$\chi$ field from $b_3(\CP^2 \times \Sph^3) = 1$ is topologically
required to complete the DFD cosmology.
\end{abstract}

\tableofcontents
\newpage

%==========================================================================
\section{Statement of the Main Theorem}
\label{sec:theorem}
%==========================================================================

\begin{theorem}[Insufficiency of $\psi$ for the CDM Transfer Function]
\label{thm:main}
Let $(\mathcal{M}^4, g_{\mu\nu})$ be a Friedmann--Lema\^{i}tre--Robertson--Walker
spacetime, and let $\psi$ be the DFD refractive scalar field defined by
the field equation
\begin{equation}
\label{eq:psi-field}
\nabla \cdot \bigl[\mu\bigl(|\nabla\psi|/a_0\bigr)\,\nabla\psi\bigr]
= -\frac{8\pi G}{c^2}\,\rho_b\,,
\end{equation}
where $\mu(x) = x/\sqrt{1+x^2}$ is the DFD interpolating function,
$a_0 \approx 1.2 \times 10^{-10}\,\mathrm{m\,s}^{-2}$ is the
MOND acceleration scale, and $\rho_b$ is the baryonic mass density.

Define the matter power spectrum achievable from $\psi$ alone as
$P_\psi(k)$, and let $P_{\mathrm{CDM}}(k)$ denote the standard
$\Lambda$CDM matter power spectrum with $\Omega_c h^2 = 0.12$.

Then there exists no solution of Eq.~\eqref{eq:psi-field}, together
with any optical magnification from the $\psi$-screen mechanism, such
that $P_\psi(k) = P_{\mathrm{CDM}}(k)$ for all $k$ in the range
$0.01 \leq k \leq 0.5\hMpc$.

In particular, $P_\psi(k)/P_b(k) \leq 30$ for all $k$, whereas
$P_{\mathrm{CDM}}(k)/P_b(k) \geq 175$ at $k = 0.1\hMpc$, where
$P_b(k)$ is the baryon-only power spectrum.
\end{theorem}

The proof proceeds through five independent arguments, each of which
is individually sufficient to establish the result. Their conjunction
provides an overdetermined proof.

%==========================================================================
\section{Definitions and Setup}
\label{sec:setup}
%==========================================================================

\subsection{The DFD $\psi$ Field}

In Density Field Dynamics, gravity arises from a variable speed of light
$c(\mathbf{x}) = c_0\,e^{-\psi(\mathbf{x})}$, where $\psi$ is determined
by the energy density of matter through Eq.~\eqref{eq:psi-field}.
The field $\psi$ corresponds to the zeroth Betti number
$b_0(\CP^2 \times \Sph^3) = 1$ of the internal manifold: it is a
scalar (0-form) sector that encodes the refractive index of the vacuum.

The Newtonian potential is related to $\psi$ by $\Phi_N = -c^2\psi/2$
in the weak-field limit, and the effective DFD potential is
\begin{equation}
\label{eq:Phi-DFD}
\Phi_{\mathrm{DFD}} = \nu\bigl(|\Phi_N|/\Phi_0\bigr)\,\Phi_N\,,
\end{equation}
where $\nu(y) = 1/\mu(x)$ evaluated at $x\mu(x) = y$, and
$\Phi_0 = a_0 \ell$ for a characteristic length $\ell$.

\subsection{The Quasi-Static Limit}

For perturbations with wavelength $\lambda \ll c/H$ (sub-horizon modes),
the time derivatives in the $\psi$ field equation are suppressed relative
to spatial gradients by a factor $(H\lambda/c)^2 \ll 1$. In this regime,
Eq.~\eqref{eq:psi-field} reduces to a spatial elliptic PDE at each
time slice:
\begin{equation}
\label{eq:elliptic}
\nabla_{\mathbf{x}} \cdot \bigl[\mu(|\nabla_{\mathbf{x}}\psi|/a_0)\,
\nabla_{\mathbf{x}}\psi\bigr] = -\frac{8\pi G}{c^2}\,\rho_b(t,\mathbf{x})\,.
\end{equation}
This is the starting point for Arguments 1 and 2.

\subsection{The CDM Transfer Function}

The CDM transfer function $T(k)$ is defined by
\begin{equation}
\label{eq:transfer}
T(k) \equiv \frac{\delta(k, z=0)}{\delta(k, z=0)|_{k\to 0}}
\cdot \frac{D(z_{\mathrm{eq}})}{D(0)}\,,
\end{equation}
normalised so that $T(k) \to 1$ for $k \to 0$. For
$\Lambda$CDM with $\Omega_c h^2 = 0.12$, the transfer function
satisfies $T(0.1\hMpc) \approx 0.5$. The critical feature is that
modes entering the horizon before matter--radiation equality
($k > k_{\mathrm{eq}} \approx 0.01\hMpc$) grow logarithmically
during radiation domination due to the CDM gravitational potential
wells, yielding a net enhancement of the matter power spectrum
$P(k) \propto T^2(k)\,k^{n_s}$ by a factor $\sim 175$ at
$k = 0.1\hMpc$ relative to a universe with baryons alone.

\subsection{What CDM Provides: The M\'esz\'aros Effect}

The mechanism by which CDM enhances the power spectrum is the
M\'esz\'aros effect: CDM perturbations $\delta_c$ decouple from
the photon--baryon plasma at $z \gg z_{\mathrm{rec}}$ and grow
(logarithmically during radiation domination, linearly during matter
domination) to establish potential wells. After recombination, baryons
fall into these pre-existing wells, boosting the baryon power spectrum.
The essential features are:
\begin{enumerate}
\item CDM does not couple to photons and therefore does not oscillate.
\item CDM perturbations grow monotonically from horizon entry onward.
\item The CDM potential wells persist through recombination.
\end{enumerate}
Any replacement for CDM must reproduce all three properties.

%==========================================================================
\section{Argument 1: Elliptic Constraint Character}
\label{sec:arg1}
%==========================================================================

\begin{lemma}[Elliptic PDEs as Instantaneous Constraints]
\label{lem:elliptic}
Let $L[u] = f$ be an elliptic PDE on a spatial domain $\Omega$ at fixed
time $t$, where $f = f(t, \mathbf{x})$ depends on time only as a parameter.
Then $u(t, \mathbf{x})$ is determined instantaneously by $f(t, \cdot)$;
that is, $u$ has no independent dynamical degrees of freedom beyond
those contained in $f$.
\end{lemma}

\begin{proof}
By the definition of ellipticity, the principal symbol
$\sigma_L(\xi) > 0$ for all $\xi \neq 0$. The Cauchy problem for
$L[u] = f$ is therefore ill-posed in the sense of Hadamard: specifying
$u$ and $\partial_t u$ on an initial surface does not determine a unique
forward evolution. Instead, given appropriate boundary conditions on
$\partial\Omega$, the elliptic regularity theorem (see Gilbarg--Trudinger,
Theorem 8.8) guarantees a unique solution $u[\cdot; f(t,\cdot)]$ that
depends continuously on $f(t, \cdot)$. The solution at time $t$ depends
only on $f$ at time $t$, with no ``memory'' of previous states.
\end{proof}

\begin{proposition}[The $\psi$ Field is a Constraint Variable]
\label{prop:psi-constraint}
In the quasi-static limit, the DFD field equation \eqref{eq:elliptic}
is a quasilinear elliptic PDE. Consequently, $\psi(t,\mathbf{x})$
is determined instantaneously by $\rho_b(t,\mathbf{x})$.
\end{proposition}

\begin{proof}
Write $F[\psi] \equiv \nabla \cdot [\mu(|\nabla\psi|/a_0)\,\nabla\psi]$.
The linearisation around a background solution $\bar{\psi}$ is
\begin{equation}
\frac{\delta F}{\delta \psi}\bigg|_{\bar{\psi}} \delta\psi
= \nabla \cdot \bigl[D_{ij}(\nabla\bar{\psi})\,\nabla_j(\delta\psi)\bigr],
\end{equation}
where the diffusion tensor is
\begin{equation}
\label{eq:diffusion-tensor}
D_{ij} = \mu(s)\,\delta_{ij} + \mu'(s)\,\frac{\partial_i\bar{\psi}\,
\partial_j\bar{\psi}}{a_0\,s}\,,
\qquad s \equiv |\nabla\bar{\psi}|/a_0\,.
\end{equation}
Since $\mu(s) > 0$ and $\mu(s) + s\mu'(s) = (1+s^2)^{-1/2} > 0$
for all $s \geq 0$, the eigenvalues of $D_{ij}$ are strictly positive.
Hence the linearised operator is uniformly elliptic for any $\bar{\psi}$,
and by Lemma~\ref{lem:elliptic}, $\delta\psi$ is instantaneously
determined by $\delta\rho_b$.

The full nonlinear operator inherits this character: by the
Leray--Schauder theorem, the quasilinear elliptic equation
$F[\psi] = -8\pi G\rho_b/c^2$ with suitable boundary conditions has a
unique solution $\psi = \psi[\rho_b]$ at each time slice.
\end{proof}

\begin{corollary}
\label{cor:no-independent-growth}
The $\psi$ field cannot grow independently of the baryon density
$\rho_b$. In particular, before recombination when $\rho_b$ is
tightly coupled to the photon--baryon plasma and undergoes acoustic
oscillations, $\psi$ must oscillate in phase with $\rho_b$. There is
no mechanism for $\psi$ to accumulate a monotonically growing
perturbation analogous to $\delta_{\mathrm{CDM}}$.
\end{corollary}

\begin{proof}
Since $\psi = \psi[\rho_b]$ instantaneously (Proposition~\ref{prop:psi-constraint}),
every property of $\psi$ is inherited from $\rho_b$. If $\rho_b$
oscillates, so does $\psi$. If $\rho_b$ has zero time-averaged
perturbation (as it does in the acoustic regime, where
$\langle\delta\rho_b\rangle_t \approx 0$ over a full oscillation cycle),
then $\langle\delta\psi\rangle_t \approx 0$ as well, up to nonlinear
corrections bounded by Argument 2 below.
\end{proof}

\begin{remark}
This is the fundamental distinction between $\psi$ and CDM. The CDM
density $\rho_c$ satisfies a \emph{hyperbolic} equation (the collisionless
Boltzmann equation), which admits freely propagating solutions. The
equation $\partial_t^2 \delta_c + 2H\partial_t\delta_c - 4\pi G\rho_c\delta_c = 0$
is a damped wave equation with a growing mode $\delta_c \propto a$
during matter domination. No analogous growing mode exists for an
elliptic constraint.
\end{remark}

%==========================================================================
\section{Argument 2: Odd Symmetry and Vanishing DC Component}
\label{sec:arg2}
%==========================================================================

\begin{definition}
\label{def:odd}
A function $f: \mathbb{R} \to \mathbb{R}$ is \emph{odd} if
$f(-x) = -f(x)$ for all $x$.
\end{definition}

\begin{lemma}[Time-Average of Odd Functions of Zero-Mean Oscillations]
\label{lem:odd-zero}
Let $f: \mathbb{R} \to \mathbb{R}$ be odd and continuous. Let
$\phi(t) = A\cos(\omega t)$ be a zero-mean oscillation. Then
\begin{equation}
\langle f(\phi) \rangle \equiv \frac{1}{T}\int_0^T f\bigl(\phi(t)\bigr)\,dt = 0\,,
\qquad T = 2\pi/\omega\,.
\end{equation}
\end{lemma}

\begin{proof}
Under the substitution $t \to T - t$, we have $\phi(T-t) = A\cos(\omega T - \omega t)
= A\cos(2\pi - \omega t) = A\cos(\omega t) = \phi(t)$ for the
fundamental period. More directly, write $\phi(t) = A\cos(\omega t)$
and substitute $u = \omega t$:
\begin{equation}
\langle f(\phi)\rangle = \frac{1}{2\pi}\int_0^{2\pi} f(A\cos u)\,du\,.
\end{equation}
Under $u \to \pi + u$, we have $\cos(\pi + u) = -\cos u$, so
$f(A\cos(\pi+u)) = f(-A\cos u) = -f(A\cos u)$ by the odd symmetry
of $f$. The integral over $[\pi, 2\pi]$ therefore cancels the integral
over $[0,\pi]$ exactly.
\end{proof}

\begin{remark}
The result extends immediately to any zero-mean periodic function
$\phi(t)$ with the symmetry $\phi(t + T/2) = -\phi(t)$, which holds
for the dominant mode of baryon acoustic oscillations.
\end{remark}

\begin{proposition}[Vanishing DC Component of the DFD Enhanced Potential]
\label{prop:odd-DFD}
Define the DFD flux function
\begin{equation}
\label{eq:flux}
f(\Phi) \equiv \nu(|\Phi|/\Phi_0)\,\Phi\,,
\end{equation}
where $\nu(y) \geq 1$ is the MOND enhancement factor. Then:
\begin{enumerate}
\item[(a)] $f$ is an odd function of $\Phi$: $f(-\Phi) = -f(\Phi)$.
\item[(b)] If $\Phi_N(t)$ oscillates with zero mean (as for the baryon
acoustic oscillation potential before recombination), then the
time-averaged DFD potential satisfies
$\langle\Phi_{\mathrm{DFD}}\rangle = \langle f(\Phi_N)\rangle = 0$.
\end{enumerate}
\end{proposition}

\begin{proof}
(a) From Eq.~\eqref{eq:flux}:
\begin{equation}
f(-\Phi) = \nu(|-\Phi|/\Phi_0)\,(-\Phi) = \nu(|\Phi|/\Phi_0)\,(-\Phi)
= -\nu(|\Phi|/\Phi_0)\,\Phi = -f(\Phi)\,.
\end{equation}

(b) This follows immediately from (a) and Lemma~\ref{lem:odd-zero}.
During the acoustic oscillation epoch, the gravitational potential
sourced by the baryon--photon plasma oscillates as
$\Phi_N(t, k) \approx A(k)\cos(c_s k\eta)$, where $c_s$ is the
sound speed and $\eta$ is conformal time. The time-average over one
oscillation period vanishes:
\begin{equation}
\langle\Phi_{\mathrm{DFD}}(t,k)\rangle
= \langle f(\Phi_N(t,k))\rangle = 0\,.
\end{equation}
\end{proof}

\begin{corollary}
\label{cor:no-potential-well}
The $\psi$-enhanced gravitational potential cannot form a static
or growing potential well during the acoustic oscillation epoch.
Any nonlinear enhancement from $\nu > 1$ is symmetric between
overdensities and underdensities, and the net time-averaged effect
on structure growth is zero.
\end{corollary}

\begin{proof}
Structure growth requires a net positive (attractive) time-averaged
potential perturbation to drive the gravitational instability:
$\ddot{\delta} + 2H\dot{\delta} = -k^2\langle\Phi\rangle/a^2$.
Since $\langle\Phi_{\mathrm{DFD}}\rangle = 0$ by
Proposition~\ref{prop:odd-DFD}, the driving term vanishes, and no
net growth accumulates over the oscillation epoch. This is in
direct contrast to CDM, where $\langle\Phi_{\mathrm{CDM}}\rangle \neq 0$
because $\delta_{\mathrm{CDM}}$ grows monotonically and is
non-oscillating.
\end{proof}

%==========================================================================
\section{Argument 3: Negligible Temporal $\psi$ Amplitude}
\label{sec:arg3}
%==========================================================================

\begin{proposition}[Amplitude Bound on the Temporal $\psi$ Sector]
\label{prop:temporal-bound}
The dust-like (temporal) branch of the $\psi$ field, as derived in
Appendix Q of DFD v3.3 (Theorem Q.7), has equation of state
$w \to 0$, $c_s^2 \to 0$ in the cosmological limit, but its energy
density satisfies
\begin{equation}
\label{eq:Omega-bound}
\Omega_\psi \equiv \frac{\rho_\psi}{\rho_{\mathrm{crit}}} < 10^{-19}\,.
\end{equation}
\end{proposition}

\begin{proof}
Theorem Q.7 of DFD v3.3 establishes the conservation law for the
temporal $\psi$ sector:
\begin{equation}
\label{eq:conservation}
a^3\,\mu(\Delta) = C_0 = \mathrm{const}\,,
\end{equation}
where $\Delta \equiv |\dot{\psi}|/(a_0 c)$ is the dimensionless
temporal gradient and $C_0$ is set by initial conditions.

The energy density of this sector is
\begin{equation}
\rho_\psi = \frac{a_0^2}{8\pi G}\,F(\Delta)\,,
\end{equation}
where $F(\Delta) = \int_0^\Delta s\mu'(s)\,ds + \Delta\mu(\Delta) - F_0$
is the Legendre transform of the MOND Lagrangian density, with
$F(\Delta) \sim \Delta^2/2$ for $\Delta \ll 1$ and
$F(\Delta) \sim \Delta$ for $\Delta \gg 1$.

From Eq.~\eqref{eq:conservation}, as $a \to \infty$ we have
$\mu(\Delta) = C_0/a^3 \to 0$, which forces $\Delta \to 0$
(since $\mu$ is monotonically increasing with $\mu(0) = 0$).
In this regime $\mu(\Delta) \approx \Delta$, so
$\Delta \approx C_0/a^3$ and
\begin{equation}
\rho_\psi \approx \frac{a_0^2}{8\pi G}\,\frac{\Delta^2}{2}
= \frac{a_0^2}{16\pi G}\,\frac{C_0^2}{a^6}\,.
\end{equation}

This falls as $a^{-6}$ --- faster than radiation ($a^{-4}$) or matter
($a^{-3}$). To estimate $C_0$, we note that the $\psi$ field is sourced
only by baryonic density fluctuations. In the early universe, the
homogeneous component of $\psi$ satisfies
$\dot{\psi} \sim \Phi_N H \sim 10^{-5} H$, giving
$\Delta_{\mathrm{ini}} \sim H/(a_0 c) \times 10^{-5}$.

At radiation--matter equality ($a_{\mathrm{eq}} \approx 3 \times 10^{-4}$),
$H_{\mathrm{eq}} \approx 10^{-14}\,\mathrm{s}^{-1}$, so
\begin{equation}
\Delta_{\mathrm{eq}} \sim \frac{10^{-5} \times 10^{-14}}{1.2\times 10^{-10}
\times 3\times 10^8} \sim 3 \times 10^{-18}\,.
\end{equation}
Since $\Delta_{\mathrm{eq}} \ll 1$, we are deep in the Newtonian regime
where $\mu(\Delta) \approx \Delta$, so
$C_0 = a_{\mathrm{eq}}^3\,\Delta_{\mathrm{eq}}
\approx (3\times 10^{-4})^3 \times 3\times 10^{-18}
\approx 8 \times 10^{-29}$.

At the present epoch ($a = 1$):
\begin{equation}
\rho_\psi(a=1) \approx \frac{a_0^2}{16\pi G}\,C_0^2
= \frac{(1.2\times 10^{-10})^2}{16\pi \times 6.67\times 10^{-11}}
\times (8\times 10^{-29})^2\,.
\end{equation}
Evaluating:
\begin{equation}
\frac{a_0^2}{16\pi G} = \frac{1.44\times 10^{-20}}{3.35\times 10^{-9}}
\approx 4.3\times 10^{-12}\,\mathrm{kg\,m}^{-3}\,,
\end{equation}
so
\begin{equation}
\rho_\psi \approx 4.3\times 10^{-12} \times 6.4\times 10^{-57}
\approx 2.8\times 10^{-68}\,\mathrm{kg\,m}^{-3}\,.
\end{equation}
Comparing to $\rho_{\mathrm{crit}} = 3H_0^2/(8\pi G)
\approx 8.5\times 10^{-27}\,\mathrm{kg\,m}^{-3}$:
\begin{equation}
\Omega_\psi = \frac{\rho_\psi}{\rho_{\mathrm{crit}}}
\approx \frac{2.8\times 10^{-68}}{8.5\times 10^{-27}}
\approx 3.3\times 10^{-42}\,.
\end{equation}
This is enormously below $10^{-19}$, confirming Eq.~\eqref{eq:Omega-bound}
with a margin of over 20 orders of magnitude.
\end{proof}

\begin{corollary}
The temporal $\psi$ sector, despite having a CDM-like equation of state
$(w \approx 0)$, contributes a fraction $\Omega_\psi/\Omega_c < 10^{-40}$
of the required CDM density. It is cosmologically irrelevant.
\end{corollary}

%==========================================================================
\section{Argument 4: Quantitative Deficit}
\label{sec:arg4}
%==========================================================================

\begin{proposition}[Quantitative Enhancement Bound]
\label{prop:quantitative}
The maximum enhancement of the matter power spectrum achievable from
the $\psi$ field alone, including both the nonlinear MOND enhancement
and the $\psi$-screen optical magnification, satisfies
\begin{equation}
\label{eq:deficit}
\frac{P_\psi(k)}{P_b(k)} \leq 30 \qquad
\text{at } k = 0.1\hMpc\,,
\end{equation}
whereas the CDM transfer function yields
\begin{equation}
\frac{P_{\mathrm{CDM}}(k)}{P_b(k)} \approx 175 \qquad
\text{at } k = 0.1\hMpc\,.
\end{equation}
The deficit factor is at least $175/30 \approx 5.8$.
\end{proposition}

\begin{proof}
We bound each contribution separately.

\textbf{(i) MOND/nonlinear enhancement.}
In the deep MOND regime ($|\nabla\Phi_N| \ll a_0$), the effective
gravitational constant is enhanced:
\begin{equation}
G_{\mathrm{eff}} = \frac{G}{\mu(|\nabla\Phi_N|/a_0)}
\approx G\,\sqrt{\frac{a_0}{|\nabla\Phi_N|/\ell}}\,.
\end{equation}
For structures at the scale $k = 0.1\hMpc$ with typical density
contrast $\delta \sim 10^{-3}$ at recombination, the Newtonian
acceleration is $g_N \sim 4\pi G\rho_b\delta/k \sim 10^{-14}\,\mathrm{m\,s}^{-2}$,
giving $g_N/a_0 \sim 10^{-4}$. In the deep MOND limit:
\begin{equation}
\frac{G_{\mathrm{eff}}}{G} \approx \frac{1}{\mu}
\approx \sqrt{\frac{a_0}{g_N}} \sim \sqrt{10^4} \sim 100\,.
\end{equation}
However, this is the enhancement of the \emph{force}, not the power
spectrum. The power spectrum enhancement scales as the square root
of the force enhancement (since $\delta \propto \Phi/c_s^2$ and
$P \propto \delta^2$), but is further suppressed by the fact that
the MOND enhancement applies only during the matter-dominated era
(not during radiation domination, when the modes of interest enter
the horizon). Accounting for the limited duration:
\begin{equation}
\text{MOND boost to } P(k): \quad \lesssim 15\times\,.
\end{equation}
This is an optimistic upper bound. The actual enhancement is
scale-independent (unlike CDM, which provides scale-dependent
enhancement through the M\'esz\'aros effect), making it impossible
to reproduce the CDM transfer function shape.

\textbf{(ii) $\psi$-screen optical magnification.}
The $\psi$-screen mechanism (Section 9 of DFD v3.3) enhances angular
diameter and luminosity distances through the refractive index
$n = e^{2\psi}$ along the line of sight. At $z \sim 1$, this provides
an effective magnification:
\begin{equation}
\frac{d_{L,\psi}}{d_{L,\mathrm{bare}}} \approx 2\,.
\end{equation}
Even if this optical effect could somehow be repurposed for structure
formation (which it cannot --- see Argument 5), it provides at most a
$2\times$ boost.

\textbf{(iii) Combined bound.}
Taking the product of both effects:
\begin{equation}
\frac{P_\psi}{P_b} \leq 15 \times 2 = 30\,.
\end{equation}

\textbf{(iv) The deficit.}
The required enhancement is $\sim 175\times$. The achievable enhancement
is $\leq 30\times$. The deficit factor is
\begin{equation}
\frac{175}{30} \approx 5.8\,.
\end{equation}
This is a \emph{lower bound} on the true deficit, because:
\begin{itemize}
\item The MOND enhancement is scale-independent, while the CDM
      transfer function is strongly scale-dependent ($T(k)$ suppresses
      small scales relative to large scales). Matching the \emph{shape}
      of $P_{\mathrm{CDM}}(k)$ is impossible with a scale-independent
      boost.
\item Argument 2 shows the time-averaged MOND enhancement vanishes
      during the oscillation epoch, so the actual MOND contribution
      to the transfer function is even smaller than the $15\times$
      bound.
\item Argument 1 shows $\psi$ cannot grow independently, further
      reducing the available enhancement window.
\end{itemize}
\end{proof}

%==========================================================================
\section{Argument 5: $\psi$-Screen Insufficiency}
\label{sec:arg5}
%==========================================================================

\begin{proposition}[$\psi$-Screen Cannot Substitute for CDM]
\label{prop:screen}
The $\psi$-screen mechanism provides an optical distance magnification
of $\sim 2\times$ at $z \sim 1$. This effect:
\begin{enumerate}
\item[(a)] is an optical effect (photon propagation through a refractive
medium), not a gravitational effect on matter clustering;
\item[(b)] does not modify the growth rate of density perturbations;
\item[(c)] cannot produce the scale-dependent transfer function shape
required to match $P_{\mathrm{CDM}}(k)$;
\item[(d)] provides a factor of $\sim 2$, whereas a factor of $\sim 175$
is required.
\end{enumerate}
\end{proposition}

\begin{proof}
\textbf{(a)} The $\psi$-screen modifies the effective metric experienced
by photons: $ds^2_{\mathrm{opt}} = n^2\,ds^2_{\mathrm{grav}}$, where
$n = e^{2\psi}$ is the refractive index. This changes the
angular diameter distance:
\begin{equation}
d_A^{\mathrm{eff}}(z) = \int_0^z \frac{c\,dz'}{H(z')\,n(\bar{\psi}(z'))}\,,
\end{equation}
but does \emph{not} modify the Poisson equation, the continuity
equation, or the Euler equation governing matter perturbation growth.
The growth factor $D(a)$ satisfies
$\ddot{D} + 2H\dot{D} - 4\pi G\rho_m D = 0$, which is independent
of the optical refractive index.

\textbf{(b)} Since the perturbation growth equations are unaffected,
the matter power spectrum $P(k,z)$ at any given redshift is identical
with or without the $\psi$-screen. The screen affects only the
\emph{mapping} between physical scales and observed angular scales,
not the amplitude or shape of $P(k)$ itself.

\textbf{(c)} The $\psi$-screen magnification is a smooth, slowly-varying
function of redshift, approximately constant across $k$-modes observed
at the same redshift. It cannot produce the strong scale dependence
of $T_{\mathrm{CDM}}(k)$, which varies by orders of magnitude between
$k = 0.01$ and $k = 0.5\hMpc$.

\textbf{(d)} Quantitatively, even if the $\psi$-screen magnification
could be applied to the power spectrum (which it cannot, by parts
(a)--(c)), the maximum boost is a factor of $\sim 2$ at $z \sim 1$,
falling to $\sim 1$ at both $z = 0$ and $z \gg 1$. This is $\sim 90\times$
too small to account for the CDM transfer function.
\end{proof}

%==========================================================================
\section{Synthesis: Completion of the Proof}
\label{sec:synthesis}
%==========================================================================

\begin{proof}[Proof of Theorem~\ref{thm:main}]
We combine the five arguments:

\textbf{Step 1} (Argument 1, \S\ref{sec:arg1}): The $\psi$ field equation
is elliptic in the quasi-static limit. Therefore $\psi$ is an instantaneous
constraint determined by $\rho_b$, with no independent dynamical degrees
of freedom. It cannot grow independently of the baryon--photon plasma.

\textbf{Step 2} (Argument 2, \S\ref{sec:arg2}): Even if the nonlinear
MOND enhancement factor $\nu > 1$ could boost the potential, the
DFD enhanced potential $\Phi_{\mathrm{DFD}} = \nu(|\Phi_N|)\Phi_N$ is
an odd function of $\Phi_N$. For oscillating baryon perturbations,
the time-averaged enhanced potential vanishes identically. No net
potential well accumulates during the acoustic epoch.

\textbf{Step 3} (Argument 3, \S\ref{sec:arg3}): The temporal $\psi$
sector (the dust branch from Appendix Q of v3.3) has $w \approx 0$ and
$c_s^2 \approx 0$, mimicking CDM kinematically. However, its amplitude
is $\Omega_\psi < 10^{-19}$ (in fact $\sim 10^{-42}$), which is
cosmologically negligible.

\textbf{Step 4} (Argument 4, \S\ref{sec:arg4}): Quantitatively, the
maximum power spectrum enhancement from the MOND nonlinearity
($\leq 15\times$) combined with the $\psi$-screen ($\leq 2\times$)
yields at most $30\times$, whereas $175\times$ is required at
$k = 0.1\hMpc$. The deficit factor is $\geq 5.8$.

\textbf{Step 5} (Argument 5, \S\ref{sec:arg5}): The $\psi$-screen is
an optical effect that modifies photon propagation, not matter
clustering. It does not alter the growth factor, cannot produce
scale-dependent enhancement, and provides only $\sim 2\times$
magnification.

Each argument independently establishes that $P_\psi(k) \neq P_{\mathrm{CDM}}(k)$.
Their conjunction provides an overdetermined proof covering the
structural (Arguments 1, 2), amplitude (Arguments 3, 4), and
categorical (Argument 5) dimensions of the failure.
\end{proof}

%==========================================================================
\section{Corollary: Topological Necessity of the $\chi$ Field}
\label{sec:corollary}
%==========================================================================

\begin{corollary}[Topological Necessity of $\chi$]
\label{cor:chi-required}
Let the DFD internal manifold be $\mathcal{K} = \CP^2 \times \Sph^3$.
Then the dark matter sector of DFD must arise from a topological sector
distinct from $\psi$. The unique candidate is the $\chi$ field
associated with the third Betti number $b_3(\CP^2 \times \Sph^3) = 1$.
\end{corollary}

\begin{proof}
The Betti numbers of $\CP^2 \times \Sph^3$ are:
\begin{equation}
\begin{aligned}
b_0 &= 1 \quad &\longrightarrow\quad &\text{scalar sector } (\psi), \\
b_1 &= 0 \quad &\longrightarrow\quad &\text{no 1-form sector}, \\
b_2 &= 1 \quad &\longrightarrow\quad &\text{gauge sector (electromagnetism)}, \\
b_3 &= 1 \quad &\longrightarrow\quad &\text{pseudoscalar sector } (\chi), \\
b_4 &= 1 \quad &\longrightarrow\quad &\text{topological density (Hodge dual of } b_0).
\end{aligned}
\end{equation}

By Theorem~\ref{thm:main}, the $b_0$ sector ($\psi$) cannot provide
the CDM transfer function. The $b_1 = 0$ sector does not exist.
The $b_2 = 1$ sector provides electromagnetism, not dark matter.
The $b_4$ sector is the Hodge dual of $b_0$ and provides no
additional dynamical degrees of freedom beyond $\psi$.

The only remaining topological sector is $b_3 = 1$, which yields
the pseudoscalar 3-form field $\chi$ on $\Sph^3$. Since DFD must
reproduce the observed matter power spectrum, and $\psi$ alone
cannot do so, $\chi$ is the unique topologically available candidate
for the dark matter role. The existence of $\chi$ is not an
\emph{ad hoc} addition but a \emph{topological consequence} of the
internal manifold $\CP^2 \times \Sph^3$.
\end{proof}

%==========================================================================
\section{Discussion: What $\chi$ Provides That $\psi$ Cannot}
\label{sec:discussion}
%==========================================================================

The five arguments of Theorem~\ref{thm:main} identify precisely
what properties a successful dark matter replacement must have, and
why the $\chi$ field from $b_3 = 1$ satisfies them while $\psi$
from $b_0 = 1$ does not.

\subsection{Independent Dynamics}

The $\chi$ field satisfies a \emph{hyperbolic} field equation
(Klein--Gordon type):
\begin{equation}
\ddot{\chi} + 3H\dot{\chi} + m_\chi^2 \chi = 0\,,
\end{equation}
with $m_\chi \approx 95.8\keV$. This is an \emph{evolution equation},
not a constraint. The $\chi$ field has two independent initial data
$(\chi_0, \dot{\chi}_0)$ per mode, and its perturbations can grow
independently of the baryon--photon plasma. Contrast this with
Argument 1: $\psi$ is determined instantaneously by $\rho_b$ and
has no independent degrees of freedom.

\subsection{Non-Oscillating Potential Wells}

Once $\chi$ becomes non-relativistic ($H < m_\chi c^2/\hbar$, which
occurs extremely early for $m_\chi \sim 100\keV$), its density
perturbations $\delta_\chi$ evolve according to the pressureless
growth equation:
\begin{equation}
\ddot{\delta}_\chi + 2H\dot{\delta}_\chi - 4\pi G\rho_\chi\delta_\chi = 0\,.
\end{equation}
The growing mode is monotonically increasing, providing the
non-oscillating potential wells that Argument 2 proves $\psi$
cannot supply.

\subsection{Correct Amplitude}

The DFD framework predicts $\Omega_\chi h^2 = (16/3)\,\Omega_b h^2
\approx 0.119$, matching the Planck-measured dark matter density.
This provides the correct amplitude for the CDM transfer function,
in contrast to Argument 3's bound $\Omega_\psi < 10^{-19}$.

\subsection{Scale-Dependent Transfer Function}

Because $\chi$ perturbations enter the horizon and begin growing at
different times depending on $k$, the resulting transfer function
$T_\chi(k)$ is naturally scale-dependent, with the characteristic
turnover at $k_{\mathrm{eq}}$ and logarithmic growth for
$k > k_{\mathrm{eq}}$. This addresses the scale-independence
failure identified in Argument 4.

\subsection{Gravitational, Not Optical}

The $\chi$ field clusters gravitationally, modifying the Poisson
equation:
\begin{equation}
\nabla^2\Phi = 4\pi G\,a^2\,(\rho_b\delta_b + \rho_\chi\delta_\chi)\,.
\end{equation}
This is a direct gravitational effect on the potential wells, not
an optical effect on photon propagation. The $\chi$ contribution
drives matter perturbation growth through the standard gravitational
instability mechanism, addressing the categorical failure identified
in Argument 5.

\subsection{Summary}

\begin{center}
\begin{tabular}{lcc}
\toprule
\textbf{Property} & \textbf{$\psi$ ($b_0$)} & \textbf{$\chi$ ($b_3$)} \\
\midrule
Field equation type & Elliptic (constraint) & Hyperbolic (evolution) \\
Independent dynamics & No & Yes \\
Potential wells persist & No ($\langle\Phi\rangle = 0$) & Yes (monotonic growth) \\
Amplitude $\Omega$ & $< 10^{-19}$ & $0.119/h^2$ \\
Scale dependence & No & Yes (M\'esz\'aros) \\
Mechanism & Optical + nonlinear & Gravitational clustering \\
\bottomrule
\end{tabular}
\end{center}

\medskip

The conclusion is unambiguous: the $\chi$ field from
$b_3(\CP^2 \times \Sph^3) = 1$ is not merely a convenient addition
to DFD, but a \emph{topological necessity} demanded by the structure
of the observed matter power spectrum. The proof presented here,
through five independent arguments, establishes that no configuration
of the $\psi$ field --- whether through nonlinear MOND enhancement,
temporal dust modes, or optical magnification --- can substitute for
the role that $\chi$ plays in DFD cosmology.

\end{document}
